\documentclass[11pt]{article}
\usepackage[T1]{fontenc}
\usepackage[a4paper,margin=1in]{geometry}
\usepackage{amsmath,amssymb,amsthm,mathtools}
\usepackage{graphicx}
\usepackage[hidelinks]{hyperref}
\usepackage{microtype}

\newtheorem{theorem}{Theorem}
\newtheorem{lemma}{Lemma}
\newtheorem{corollary}{Corollary}
\newcommand{\C}{\mathbb C}
\newcommand{\N}{\mathbb N}
\newcommand{\sap}{\sigma_{\mathrm{ap}}}
\newcommand{\sbap}{\sigma_{\mathrm{bap}}}
\newcommand{\spoint}{\sigma_{\mathrm p}}
\newcommand{\saseq}{\sigma_{\mathrm{ap}}^{\mathrm{seq}}}
\newcommand{\sbseq}{\sigma_{\mathrm{bap}}^{\mathrm{seq}}}

\setlength{\parindent}{0pt}
\setlength{\parskip}{0.55em}

\title{A Non-Banach Class with Complete Collapse of Approximate Spectra}
\author{Run \texttt{fa\_banach\_001} --- lane 3}
\date{19 July 2026}

\begin{document}
\maketitle

\begin{center}
\fbox{\parbox{0.9\textwidth}{
\textbf{Status: claimed full answer, likely valid.}
If a complex vector space carries its strongest Hausdorff locally convex
topology, then for every linear operator all four approximate point spectra
in Remark~5.17(a) equal the point spectrum.  The standard space
$\varphi=\bigoplus_{n\geq1}\C$ is complete, infinite-dimensional, and
nonnormable, so this is genuinely beyond the Banach setting.
}}
\end{center}

\section{Source question}

Karsten Kruse's \emph{Spectral theory for semigroups on locally convex spaces},
arXiv:2509.11931v2, Remark~5.17(a), asks for useful sufficient conditions,
apart from the Banach-space case, which ensure
\[
 \sap(B)=\sbap(B)
 \qquad\hbox{or}\qquad
 \saseq(B)=\sbseq(B)
\]
for every closed linear operator $(B,D(B))$ on a Hausdorff locally convex
space $X$~\cite{Kruse}.

\begin{figure}[ht]
  \centering
  \includegraphics[width=0.96\textwidth]{figures/open_problem_crop.png}
  \caption{Remark~5.17(a) in arXiv:2509.11931v2, PDF page 34.}
\end{figure}

\section{Definitions and answer}

The \emph{strongest Hausdorff locally convex topology} on a complex vector
space $X$ is the topology generated by all seminorms on $X$.  Equivalently,
every linear map from $X$ into any locally convex space is continuous.

For a linear operator $(B,D(B))$ and $\lambda\in\C$, the source defines
$\lambda\in\sap(B)$ when there is a net $(x_i)\subseteq D(B)$ which does not
converge to zero but satisfies $(B-\lambda)x_i\to0$.  The symbol $\sbap$
requires the witnessing net to be bounded, and the superscript
$\mathrm{seq}$ restricts to sequences.  The point spectrum $\spoint(B)$
consists of the eigenvalues.

\section{Proof intuition}

An approximate eigenvalue outside the point spectrum is exactly a failure of
continuity of the algebraic inverse
\[
 (\lambda-B)^{-1}\colon \operatorname{ran}(\lambda-B)\longrightarrow X.
\]
On a space with the strongest locally convex topology, every subspace again
has the strongest locally convex topology.  Therefore every linear map out of
the range---in particular this inverse---is automatically continuous.  No
nonzero approximate eigen-net can exist unless a genuine eigenvector already
exists.  This argument does not need boundedness, sequentiality, or closedness
of the operator.

\section{The full collapse theorem}

\begin{lemma}[heredity of the strongest topology]\label{lem:heredity}
Let $X$ carry its strongest Hausdorff locally convex topology and let
$Y\subseteq X$ be a linear subspace.  Then the topology induced on $Y$ is the
strongest Hausdorff locally convex topology on $Y$.  Consequently, every
linear map from $Y$ to a Hausdorff locally convex space is continuous.
\end{lemma}

\begin{proof}
Let $q$ be any seminorm on $Y$.  Choose an algebraic complement $Z$ with
$X=Y\oplus Z$, and let $P\colon X\to Y$ be the associated algebraic
projection.  Then $q\circ P$ is a seminorm on $X$, hence is continuous by the
definition of the topology of $X$.  Its restriction to $Y$ is $q$, so every
seminorm on $Y$ is continuous for the induced topology.  This is precisely the
strongest locally convex topology on $Y$.

If $L\colon Y\to F$ is linear and $p$ is a continuous seminorm on $F$, then
$p\circ L$ is a seminorm on $Y$ and is therefore continuous.  Hence $L$ is
continuous.
\end{proof}

\begin{theorem}\label{thm:collapse}
Let $X$ be a complex vector space with its strongest Hausdorff locally convex
topology.  For every linear operator $(B,D(B))$ on $X$---not necessarily
closed or densely defined---one has
\[
 \boxed{
 \spoint(B)=\sap(B)=\saseq(B)=\sbap(B)=\sbseq(B).}
\]
In particular, both equalities requested in Remark~5.17(a) hold for every
closed operator on $X$.
\end{theorem}

\begin{proof}
Fix $\lambda\notin\spoint(B)$.  Then
$T:=\lambda-B\colon D(B)\to X$ is injective.  Put
$R:=\operatorname{ran}T\subseteq X$.  The algebraic inverse
\[
 S:=T^{-1}\colon R\longrightarrow X
\]
is linear.  By Lemma~\ref{lem:heredity}, $R$ has the strongest locally convex
topology, and hence $S$ is continuous.

If a net $(x_i)\subseteq D(B)$ satisfies $(B-\lambda)x_i\to0$, then
\[
 x_i=S(Tx_i)=-S((B-\lambda)x_i)\longrightarrow0.
\]
Thus $\lambda\notin\sap(B)$, and therefore it belongs to none of the three
smaller approximate spectra.  This proves that all four approximate point
spectra are contained in $\spoint(B)$.

Conversely, if $Bx=\lambda x$ for some $x\ne0$, the constant net and constant
sequence with value $x$ are bounded, do not tend to zero, and have zero
residual.  Hence $\lambda$ belongs to all four approximate point spectra.
\end{proof}

\begin{corollary}[a concrete complete non-Banach class]\label{cor:phi}
Let
\[
 \varphi=\C^{(\N)}=\bigoplus_{n\geq1}\C
\]
carry its locally convex direct-sum topology, equivalently its strongest
locally convex topology.  Then $\varphi$ is complete, infinite-dimensional,
and nonnormable, and the five spectra in Theorem~\ref{thm:collapse} coincide
for every linear operator on $\varphi$.
\end{corollary}

\begin{proof}
The spectral assertion is Theorem~\ref{thm:collapse}.  The space $\varphi$ is
the strict locally convex inductive limit of the finite-dimensional spaces
$\C^n$.  It is the standard complete Montel $(DF)$-space with the strongest
locally convex topology; see, for example, K\k{a}kol--Saxon~\cite{KakolSaxon}.
It is nonnormable: an infinite-dimensional normed space admits discontinuous
algebraic linear functionals, whereas every linear functional on $\varphi$ is
continuous by definition of its topology.
\end{proof}

Thus the sufficient condition is compatible with completeness and is not a
disguised Banach assumption.  The same proof works in arbitrary Hamel
dimension for the strongest locally convex topology; $\varphi$ is singled out
because it is a standard sequential complete nonnormable example relevant to
the source's semigroup framework.

\section{Sharp contrast with Montel and nuclear hypotheses}

The strongest-topology hypothesis cannot be replaced merely by familiar
compactness or regularity assumptions.  Let
\[
 s=\left\{x=(x_n):p_k(x):=\sup_{n\ge1}n^k|x_n|<\infty
       \text{ for every }k\in\N_0\right\}
\]
be the nuclear Fr\'echet--Montel space of rapidly decreasing sequences, and
define $(Bx)_n=e^{-n}x_n$.  Then $B$ is continuous, closed, and injective, but
\[
 0\in\saseq(B)\setminus\sbap(B).
\]
Indeed, the unit vectors $e_m$ do not tend to zero since $p_0(e_m)=1$, while
$p_k(Be_m)=m^ke^{-m}\to0$ for every $k$.  Conversely, if $(x_i)$ is bounded
and $Bx_i\to0$, then for fixed $k$ and
$C=\sup_i p_{k+1}(x_i)$,
\[
 \sup_{n>N}n^k|x_{i,n}|\le C/N,
\]
whereas the finitely many coordinates $n\le N$ tend to zero.  Hence
$p_k(x_i)\to0$ for every $k$, so $x_i\to0$.  This rules out completeness,
metrizability, nuclearity, Fr\'echet, Schwartz, or Montel hypotheses---even in
combination---as sufficient conditions by themselves.

\section{Scope, verification, and novelty check}

Theorem~\ref{thm:collapse} fully answers the source's request for a useful
sufficient condition beyond Banach spaces; it does not claim to characterize
all locally convex spaces having the property.  The theorem is stronger than
requested because it treats arbitrary operators and both net and sequential
versions simultaneously.

The proof was checked directly against the definitions in the current June
2026 arXiv v2.  Its only structural step, Lemma~\ref{lem:heredity}, is proved
from an algebraic complement and the defining family of all seminorms.  A
bounded novelty search through 19 July 2026 covered the run registry, the
local arXiv source corpus, the current source paper and its references, arXiv,
and exact web phrases combining ``approximate point spectrum'' with
``strongest/finest locally convex topology'' and $\varphi$.  It found the
standard topology literature and unrelated approximate-spectrum literature,
but no prior statement of Theorem~\ref{thm:collapse} or its application to
Remark~5.17(a).  The topology facts are classical; novelty of the spectral
observation remains subject to expert review.

\begin{thebibliography}{9}
\bibitem{Kruse}
K. Kruse, \emph{Spectral theory for semigroups on locally convex spaces},
J. Math. Anal. Appl. \textbf{564} (2026), Article 130889;
arXiv:2509.11931v2.

\bibitem{KakolSaxon}
J. K\k{a}kol and S. A. Saxon,
\emph{Montel $(DF)$-spaces, sequential $(LM)$-spaces and the strongest locally
convex topology}, J. London Math. Soc. (2) \textbf{66} (2002), 388--406,
doi:10.1112/S0024610702003459.
\end{thebibliography}

\end{document}
